\chapter*{Statement of Originality}
\addcontentsline{toc}{chapter}{Statement of Originality}
\vspace{-0.5cm}

The content in this thesis is an account of the work I did between 31 January 2022 and 23 February 2026, while enrolled as a PhD candidate at the Research School of Astronomy and Astrophysics, The Australian National University. Except where explicitly acknowledged, the material is my own original work.

\textbf{Chapters 1} (Introduction) and \textbf{5} (Conclusion) are original material written specifically for this thesis. In Chapter 1, a small number of schematics have been redrawn by me, based on material that appear in textbooks, lecture notes, or review articles; in each case the source of inspiration is clearly credited. Where discussions retread established arguments, appropriate references have been provided.

\textbf{Chapters 2} and \textbf{3} are lightly adapted versions of the following publications:
\begin{description}
    \item \underline{{Kriel, N.}}, Beattie, J. R., Federrath, C., Krumholz, M. R., \& Hew, J., 2025, Fundamental MHD scales -- II: the kinematic phase of the supersonic small-scale dynamo,MNRAS, 537(3), 2602–2629,doi:\href{https://doi.org/10.1093/mnras/staf188}{10.1093/mnras/staf188}
    
    \item \underline{{Kriel, N.}}, Beattie, J. R., Krumholz, M. R., Schober, J., \& Armstrong, P. J., The universal growth of magnetic energy during the nonlinear phase of subsonic and supersonic small-scale dynamos, Phys. Rev. E, accepted
\end{description}
For these two studies, I led in all aspects: I ran the simulations, carried out the data analysis, developed the theoretical interpretation, drafted the manuscripts, and coordinated the response with referees. My co-authors contributed through scientific discussion, development of ideas, and detailed feedback on the manuscripts.

\clearpage
\textbf{Chapter 4} is original work based on a study currently in preparation:
\begin{description}
    \item \underline{{Kriel, N.}}, Cole-Kodikara, E., Wibking, B., Krumholz, M. R., He, C.-C., Sarkar, A., \& Li, P. S., Magnetising the \textsc{Quokka} code with fast, second-order accurate, error-correcting schemes for magnetohydrodynamics on GPUs
\end{description}
This chapter describes the development of an MHD module for the open-source code \textsc{Quokka}. Elizabeth Cole-Kodikara performed the strong and weak scaling analysis. Mark R. Krumholz contributed draft text for parts of the manuscript, which I edited and integrated. The broader \textsc{Quokka} collaboration, including Benjamin Wibking and Chong-Chong He, provided helpful discussion and feedback throughout development, including review of code changes and contributions to testing infrastructure. All other aspects of this chapter, including the design and implementation of the MHD infrastructure, the correctness tests, the figures and schematics, and the manuscript itself, are my own.

\vspace{2 cm}
\begin{flushright}
  \noindent\rule{5.0cm}{0.4 pt}\\
  Neco Kriel\\
  Bonn, Germany, August 2026
\end{flushright}